\documentclass[11pt,a4paper]{article}

% =====================================================================
%  Statistiques & Probabilités — Synthèse complète (Terminale MC)
%  Version LaTeX du document interactif « synthese-stats-probas ».
%  Compilation : pdflatex (2 passes pour la table des matières).
% =====================================================================

\usepackage[utf8]{inputenc}
\usepackage[T1]{fontenc}
\usepackage{lmodern}
\usepackage[french]{babel}
\usepackage{amsmath,amssymb}
\usepackage[a4paper,margin=2.1cm,headheight=14pt]{geometry}
\usepackage{graphicx}
\usepackage{booktabs}
\usepackage{array}
\usepackage{tabularx}
\usepackage{enumitem}
\usepackage{xcolor}
\usepackage[most]{tcolorbox}
\usepackage{tikz}
\usetikzlibrary{arrows.meta,calc,positioning,babel}
\usepackage{titlesec}
\usepackage{microtype}
\usepackage{hyperref}

% ----------------------------- Couleurs ------------------------------
\definecolor{spInk}{HTML}{1B2A41}
\definecolor{spSoft}{HTML}{4B5B76}
\definecolor{spLink}{HTML}{1C64B0}
\definecolor{spBlue}{HTML}{4A90D9}
\definecolor{spBlueDark}{HTML}{1C64B0}
\definecolor{spGreen}{HTML}{3DA95F}
\definecolor{spGreenDark}{HTML}{1F7A3A}
\definecolor{spOrange}{HTML}{B5651D}
\definecolor{spPurple}{HTML}{7C4DBE}
\definecolor{spAmberBg}{HTML}{FCF3E3}

\hypersetup{colorlinks=true, linkcolor=spLink, urlcolor=spLink}

% --------------------------- Titres de sections ----------------------
\titleformat{\section}
{\Large\bfseries\sffamily\color{spInk}}{\thesection}{0.8em}{}
\titleformat{\subsection}
{\large\bfseries\sffamily\color{spInk}}{\thesubsection}{0.7em}{}
\titleformat{\subsubsection}
{\normalsize\bfseries\sffamily\color{spInk}}{\thesubsubsection}{0.6em}{}

% ------------------------------ Boîtes -------------------------------
\tcbset{
	fichebase/.style={
		enhanced, breakable, boxrule=0.8pt, arc=2mm,
		left=3mm, right=3mm, top=2mm, bottom=2mm,
		fonttitle=\bfseries\sffamily, coltitle=white,
		attach boxed title to top left={yshift=-2mm, xshift=3mm},
		boxed title style={arc=1.2mm, boxrule=0pt, left=2mm, right=2mm,
			top=0.8mm, bottom=0.8mm}
	}
}
\newtcolorbox{fiche}[1]{fichebase,
	colback=spBlue!5, colframe=spBlueDark!75, colbacktitle=spBlueDark, title={#1}}
\newtcolorbox{ficheoutils}[1]{fichebase,
	colback=spGreen!6, colframe=spGreenDark!80, colbacktitle=spGreenDark, title={#1}}
\newtcolorbox{ficheprobas}[1]{fichebase,
	colback=spPurple!6, colframe=spPurple!80!black, colbacktitle=spPurple, title={#1}}
\newtcolorbox{formule}{colback=white, colframe=spBlue!45, boxrule=0.6pt,
	arc=1.5mm, left=3mm, right=3mm, top=1.5mm, bottom=1.5mm, breakable,
	before skip=4pt, after skip=6pt}
\newtcolorbox{notebox}{colback=spAmberBg, colframe=spOrange!85!black,
	boxrule=0.7pt, arc=1.5mm, left=3mm, right=3mm, top=2mm, bottom=2mm,
	breakable, before skip=4pt, after skip=6pt}

% ---------------------------- Styles TikZ ----------------------------
\tikzset{
	ndbase/.style={rounded corners=3pt, draw, line width=0.5pt, align=center},
	ndRoot/.style={ndbase, fill=spInk!88, text=white, font=\bfseries\scriptsize\sffamily, text width=2.5cm},
	ndStatH/.style={ndbase, fill=spBlue!32, draw=spBlueDark, font=\bfseries\scriptsize\sffamily, text width=1.6cm},
	ndStat/.style={ndbase, fill=spBlue!12, draw=spBlueDark!75, font=\scriptsize\sffamily, text width=1.7cm},
	ndProbH/.style={ndbase, fill=spPurple!30, draw=spPurple!80!black, font=\bfseries\scriptsize\sffamily, text width=1.6cm},
	ndProb/.style={ndbase, fill=spPurple!10, draw=spPurple!70!black, font=\scriptsize\sffamily, text width=1.7cm},
	ndSampH/.style={ndbase, fill=spGreen!30, draw=spGreenDark, font=\bfseries\scriptsize\sffamily, text width=1.9cm},
	ndIfH/.style={ndbase, fill=spOrange!28, draw=spOrange!80!black, font=\bfseries\scriptsize\sffamily, text width=1.8cm},
	ndIf/.style={ndbase, fill=spOrange!10, draw=spOrange!70!black, font=\scriptsize\sffamily, text width=1.9cm},
	ndIcH/.style={ndbase, fill=spGreen!28, draw=spGreenDark, font=\bfseries\scriptsize\sffamily, text width=1.8cm},
	ndIc/.style={ndbase, fill=spGreen!10, draw=spGreenDark!75, font=\scriptsize\sffamily, text width=1.9cm},
	gstat/.style={ndbase, fill=spBlue!12, draw=spBlueDark!75, font=\small\sffamily, inner sep=3.5pt},
	gtool/.style={ndbase, fill=spGreen!12, draw=spGreenDark!80, font=\small\sffamily, inner sep=3.5pt},
	gprob/.style={ndbase, fill=spPurple!10, draw=spPurple!70!black, font=\small\sffamily, inner sep=3.5pt},
	glabel/.style={font=\tiny\sffamily, fill=white, fill opacity=0.75, text opacity=1, inner sep=1pt}
}

% ------------------------------ Raccourcis ---------------------------
\newcommand{\IF}{\textit{IF}}
\newcommand{\IC}{\textit{IC}}

% =====================================================================
\title{{\sffamily\bfseries Statistiques \& Probabilités}\\[4pt]
	{\large\sffamily Synthèse complète — cours, tableaux, carte mentale,\\
		graphe de connaissances, simulateur et quiz}\\[10pt]
	{\small Terminale générale · Option Mathématiques complémentaires}}
\author{}
\date{}

\begin{document}
	\maketitle
	\vspace{-1.5em}
	\begin{center}
		\begin{minipage}{0.86\textwidth}
			\small\itshape Le même contenu que la version interactive, organisé en six volets :
			fiches de cours, tableaux de synthèse, carte mentale, graphe de connaissances,
			simulateur retranscrit en exemples commentés, et quiz d'auto-évaluation avec corrigé.
			Chaque notion fait l'objet d'une fiche détaillée avec définitions, formules,
			exemples et pièges à éviter.
		\end{minipage}
	\end{center}
	
	\tableofcontents
	\newpage
	
	% =====================================================================
	\section{Le cours en fiches}
	% =====================================================================
	
	Définitions, interprétations et illustrations. Le fil rouge du cours est un
	\textbf{forum de jeu} qui compte $10\,000$ membres ; on y prélève un échantillon
	de $500$ membres pour étudier la faction qu'ils choisissent (Alliance ou Horde),
	et l'ouverture d'une \emph{loot box} pour les probabilités.
	
	\subsection{Vue d'ensemble : statistiques, échantillonnage et probabilités}
	
	\begin{fiche}{Synthèse : statistiques, échantillonnage et probabilités}
		Cette synthèse relie les \textbf{trois grandes familles} de notions au programme
		de Terminale (option Mathématiques Complémentaires) :
		\begin{itemize}
			\item \textbf{Statistiques descriptives} : décrire une population à partir de
			données observées (proportions, fréquences) ;
			\item \textbf{Échantillonnage} : passer d'un échantillon à la population — ou
			l'inverse — avec un contrôle du risque d'erreur (\IF{}, \IC{}) ;
			\item \textbf{Probabilités} : modéliser l'incertitude (univers, événements,
			probabilité conditionnelle).
		\end{itemize}
		Le fil conducteur : une \textbf{fréquence observée} $f$ sur un échantillon permet
		de \textbf{vérifier} une proportion $p$ (si elle est connue) ou de
		l'\textbf{estimer} (si elle est inconnue).
	\end{fiche}
	
	\subsection{Statistiques descriptives — ce que j'observe}
	
	Les statistiques descriptives organisent les \textbf{données observées}.
	
	\begin{fiche}{Population $\left(\mathcal{P}\right)$}
		La \textbf{population} est l'ensemble de tous les individus concernés par
		l'étude. On note :
		\begin{itemize}
			\item $N$ : la \textbf{taille} de la population (nombre total d'individus) ;
			\item $p$ : la \textbf{proportion théorique} d'une modalité dans cette population.
		\end{itemize}
		\begin{formule}
			\[ p = \dfrac{N_A}{N} \]
		\end{formule}
		où $N_A$ est le nombre d'individus de la population qui possèdent la modalité $A$.
		\begin{notebox}
			\textbf{Remarque :} en pratique, $p$ est souvent \emph{inconnue} — c'est ce que
			l'on cherche à estimer via un sondage. $N$ est en revanche généralement connu,
			au moins en ordre de grandeur.
		\end{notebox}
		\emph{Exemple :} les $10\,000$ membres du forum ; une population de
		$10\,000$ membres dont $3\,000$ ont choisi l'Alliance donne
		$N = 10\,000$, $N_A = 3\,000$, donc $p = \dfrac{3\,000}{10\,000} = 0{,}3$.
	\end{fiche}
	
	\begin{fiche}{Individu}
		Un \textbf{individu} est un élément de la population. C'est l'objet de base sur
		lequel on observe le caractère étudié.
		
		Chaque individu porte une \textbf{modalité} du caractère étudié. Par exemple :
		\begin{itemize}
			\item le joueur « DarkVador59 » est un individu ;
			\item sa faction (Alliance ou Horde) est sa modalité.
		\end{itemize}
	\end{fiche}
	
	\begin{fiche}{Échantillon $\left(n\right)$}
		Un \textbf{échantillon} est un sous-ensemble de la population, obtenu par tirage
		au sort (idéalement aléatoire, avec ou sans remise). On note $n$ sa
		\textbf{taille} : le nombre d'individus prélevés.
		
		L'échantillon sert à :
		\begin{itemize}
			\item \textbf{observer} une fréquence $f$ d'une modalité ;
			\item \textbf{vérifier} une hypothèse sur $p$ (Intervalle de Fluctuation) ;
			\item \textbf{estimer} $p$ quand elle est inconnue (Intervalle de Confiance).
		\end{itemize}
		\begin{notebox}
			\textbf{Important :} $n$ désigne \emph{toujours} la taille de l'échantillon.
			Pour un effectif de modalité dans l'échantillon, on écrit $n_A$ (avec un indice).
		\end{notebox}
	\end{fiche}
	
	\begin{fiche}{Caractère / Modalité}
		Le \textbf{caractère} est ce que l'on étudie sur chaque individu. Une
		\textbf{modalité} est une valeur possible du caractère.
		\begin{itemize}
			\item \textbf{Caractère qualitatif} : valeurs non numériques
			(ex. : la faction choisie $\rightarrow$ Alliance, Horde).
			\item \textbf{Caractère quantitatif} : valeurs numériques
			(ex. : l'âge, le nombre d'heures de jeu).
		\end{itemize}
		Une \emph{modalité} est une catégorie du caractère qualitatif (ex. : « Alliance »).
	\end{fiche}
	
	\begin{fiche}{Effectifs : $N_A$ et $n_A$}
		L'\textbf{effectif} d'une modalité $A$ est le \textbf{nombre d'individus} qui la
		portent :
		\begin{itemize}
			\item $N_A$ : dans la \textbf{population} ;
			\item $n_A$ : dans l'\textbf{échantillon}.
		\end{itemize}
		\begin{formule}
			\[ N_A \text{ (population)} \qquad\qquad n_A \text{ (échantillon)} \]
		\end{formule}
		\emph{Exemple :} dans le forum, $N_A = 3\,000$ membres jouent Alliance ; dans
		l'échantillon de $n = 500$ membres, $n_A = 160$ jouent Alliance.
		\begin{notebox}
			\textbf{Piège de notation :} l'indice est obligatoire ! Sans indice, $N$ et $n$
			désignent des \emph{tailles} (totaux), jamais des effectifs de modalité.
		\end{notebox}
	\end{fiche}
	
	\begin{fiche}{Fréquence $f = n_A / n$ — et proportion $p$}
		La \textbf{fréquence observée} $f$ d'une modalité $A$ dans un échantillon est :
		\begin{formule}
			\[ f = \dfrac{n_A}{n} \]
		\end{formule}
		où :
		\begin{itemize}
			\item $n_A$ est le nombre d'individus de l'échantillon portant la modalité $A$ ;
			\item $n$ est la taille totale de l'échantillon.
		\end{itemize}
		La fréquence $f$ est un nombre entre 0 et 1. C'est une \textbf{observation} sur
		l'échantillon, à ne pas confondre avec $p$, la proportion théorique dans la
		population : théorie contre observation.
		\begin{notebox}
			\textbf{Piège fréquent :} on n'écrit \emph{jamais} $f = \dfrac{n}{N}$ ! La lettre
			$n$ désigne la taille de l'échantillon, pas un effectif de modalité. Pour un
			effectif de modalité, on utilise toujours un indice : $n_A$, $n_i$, $N_A$, \ldots
		\end{notebox}
	\end{fiche}
	
	\subsubsection{Boîte à notations — la nomenclature du cours}
	
	\begin{notebox}
		\textbf{Nomenclature complète :}
		\begin{itemize}
			\item $N$ : \textbf{taille de la population} (nombre total d'individus de la
			population) ;
			\item $n$ : \textbf{taille de l'échantillon} (nombre total d'individus prélevés) ;
			\item pour une modalité $A$ — la lettre $A$ est une \emph{étiquette}
			(« aime le fromage », « vote oui »), pas un nombre — on ajoute un indice :
			\begin{itemize}
				\item $N_A$ : nombre d'individus \textbf{de la population} portant $A$ ;
				\item $n_A$ : nombre d'individus \textbf{de l'échantillon} portant $A$ ;
			\end{itemize}
			\item d'où les deux quotients à connaître par cœur :
			\[ p = \dfrac{N_A}{N} \quad\text{(proportion, population)} \qquad\qquad
			f = \dfrac{n_A}{n} \quad\text{(fréquence, échantillon)} \]
			\item \textbf{Exemple complet.} Population : $N = 10\,000$,
			$N_A = 3\,000$ (Alliance), donc $p = \dfrac{3\,000}{10\,000} = 0{,}3$.
			Échantillon : $n = 500$, $n_A = 160$, donc
			$f = \dfrac{160}{500} = 0{,}32$.
		\end{itemize}
	\end{notebox}
	
	\begin{tcolorbox}[colback=white, colframe=spOrange!90!black, boxrule=0.9pt,
		arc=2mm, left=3mm, right=3mm, top=2mm, bottom=2mm,
		title={\sffamily\bfseries À retenir}, fonttitle=\bfseries,
		coltitle=white, colbacktitle=spOrange!85!black]
		On n'écrit jamais $\dfrac{n}{N}$ pour désigner la proportion d'une modalité !
		Ce quotient compare la taille de l'échantillon à celle de la population : ce
		n'est pas une proportion. Tout effectif de modalité porte un indice :
		$n_A$, $n_i$, $N_A$, \ldots
	\end{tcolorbox}
	
	\begin{fiche}{Branche statistiques (vue d'ensemble)}
		La statistique descriptive organise et résume les données observées sur une
		population ou un échantillon. On y trouve : la \textbf{population} $\left(N\right)$,
		l'\textbf{échantillon} $\left(n\right)$, les \textbf{individus}, les
		\textbf{caractères} et \textbf{modalités}, et les \textbf{fréquences} $f$.
		\begin{notebox}
			\textbf{Attention :} ne pas confondre $n$ (taille de l'échantillon) et $n_A$
			(effectif d'une modalité dans l'échantillon).
		\end{notebox}
	\end{fiche}
	
	\subsection{Échantillonnage — le duel IF contre IC}
	\label{sec:echantillonnage}
	
	L'échantillonnage est le \textbf{pont} entre les statistiques (ce que je vois :
	$f$) et les probabilités (ce que je modélise : $p$). Deux outils, deux sens de
	lecture — c'est le fameux « chassé-croisé » :
	
	\begin{figure}[htbp]
		\centering
		\begin{tikzpicture}[>=Stealth, node distance=1cm,
			boite/.style={rounded corners=10pt, draw, line width=1.2pt, align=center,
				inner sep=7pt, font=\small}]
			\node[boite, draw=spBlue, text width=4.1cm] (pop)
			{{\bfseries\color{spBlueDark} Population}\\[2pt]
				$N$ individus $\cdot$ proportion $p$\\[1pt]
				{\itshape\color{spSoft} ce que je modélise}};
			\node[boite, draw=spGreen, text width=4.1cm, right=5.6cm of pop] (ech)
			{{\bfseries\color{spGreenDark} Échantillon}\\[2pt]
				$n$ individus $\cdot$ fréquence $f$\\[1pt]
				{\itshape\color{spSoft} ce que j'observe}};
			\draw[->, line width=1.4pt, spOrange]
			(pop.north) to[bend left=28]
			node[above=3pt, align=center, font=\small]
			{{\bfseries\color{spOrange} IF $\cdot$ vérifier}\\
				{\footnotesize $p$ est connu $\rightarrow$ $f$ est-il normal ?}}
			(ech.north);
			\draw[->, line width=1.4pt, spGreenDark]
			(ech.south) to[bend right=28]
			node[below=5pt, align=center, font=\small]
			{{\bfseries\color{spGreenDark} IC $\cdot$ estimer}\\
				{\footnotesize $f$ est observé $\rightarrow$ où se cache $p$ ?}}
			(pop.south);
		\end{tikzpicture}
		\caption{Le « chassé-croisé » de l'échantillonnage : l'IF va de la population
			vers l'échantillon pour vérifier ; l'IC va de l'échantillon vers la population
			pour estimer.}
	\end{figure}
	
	\begin{ficheoutils}{Échantillonnage (vue d'ensemble)}
		L'\textbf{échantillonnage} regroupe les méthodes qui consistent à prélever un
		échantillon dans une population et à en tirer des conclusions sur cette
		population. C'est le \textbf{pont} entre :
		\begin{itemize}
			\item les \textbf{statistiques descriptives} (ce que je vois sur l'échantillon) ;
			\item les \textbf{probabilités} (le modèle théorique qui décrit la population).
		\end{itemize}
		Deux outils complémentaires :
		\begin{itemize}
			\item l'\textbf{Intervalle de Fluctuation} (\IF{}) : de $p$ vers $f$ —
			\emph{vérifier} ;
			\item l'\textbf{Intervalle de Confiance} (\IC{}) : de $f$ vers $p$ —
			\emph{estimer}.
		\end{itemize}
	\end{ficheoutils}
	
	\subsubsection{Intervalle de Fluctuation (IF)}
	
	\begin{ficheoutils}{Intervalle de Fluctuation (IF)}
		Pour un échantillon de taille $n$, l'\textbf{intervalle de fluctuation} au seuil
		de 95~\% est l'intervalle centré sur $p$ qui contient la fréquence observée $f$
		avec une probabilité d'au moins $0{,}95$.
		\begin{formule}
			\[ \IF{} = \left[p - \dfrac{1}{\sqrt{n}}\ ;\ p + \dfrac{1}{\sqrt{n}}\right] \]
		\end{formule}
		\textbf{Usage :} décider si un échantillon observé est « normal » (compatible
		avec $p$) ou « anormal » (incompatible, donc suspect).
		\begin{notebox}
			\textbf{Conditions d'application :} $n \geq 25$ et $0{,}2 \leq p \leq 0{,}8$.
			Le seuil de 95~\% correspond à un risque d'erreur de 5~\% quand on rejette
			l'hypothèse.
		\end{notebox}
	\end{ficheoutils}
	
	\begin{ficheoutils}{IF — Quand l'utiliser ?}
		On utilise l'\textbf{intervalle de fluctuation} quand la proportion $p$ de la
		population est \textbf{connue} (ou supposée connue) et que l'on veut tester si
		un échantillon observé est cohérent avec cette valeur.
		
		\emph{Exemples de situations :}
		\begin{itemize}
			\item vérifier qu'un dé n'est pas truqué ($p = \dfrac{1}{6}$ pour chaque face) ;
			\item tester si un candidat a triché à un examen ;
			\item vérifier la conformité d'un lot de production ($p$ = taux de défaut attendu).
		\end{itemize}
		On part de la \textbf{théorie} (population) et on regarde si l'
		\textbf{observation} (échantillon) est compatible.
	\end{ficheoutils}
	
	\begin{ficheoutils}{IF — Pourquoi l'utiliser ?}
		L'IF sert à \textbf{prendre une décision} :
		\begin{itemize}
			\item \textbf{Valider} une hypothèse : si $f \in \IF{}$, l'écart entre $f$ et
			$p$ est compatible avec les fluctuations d'échantillonnage ;
			\item \textbf{Rejeter} une hypothèse : si $f \notin \IF{}$, l'écart est trop
			important pour être dû au hasard — on rejette l'hypothèse, avec un risque
			d'erreur de 5~\%.
		\end{itemize}
		C'est un \textbf{test statistique} simplifié, basé sur l'intervalle de
		fluctuation asymptotique.
	\end{ficheoutils}
	
	\begin{ficheoutils}{IF — Interprétation}
		Après calcul de l'intervalle de fluctuation
		$\IF{} = \left[p - \dfrac{1}{\sqrt{n}}\ ;\ p + \dfrac{1}{\sqrt{n}}\right]$ :
		\begin{itemize}
			\item si la fréquence observée $f$ appartient à $\IF{}$ : l'échantillon est
			considéré comme \textbf{normal}. La différence entre $f$ et $p$ s'explique
			par le hasard de l'échantillonnage ;
			\item si $f$ n'appartient pas à $\IF{}$ : l'échantillon est
			\textbf{anormal}. On \textbf{rejette l'hypothèse} selon laquelle la
			proportion vaut $p$, au risque (5~\%) de se tromper.
		\end{itemize}
		\begin{notebox}
			\textbf{Attention :} « rejeter » ne veut pas dire « prouver que c'est faux ».
			C'est une décision statistique, prise avec un risque d'erreur contrôlé.
		\end{notebox}
	\end{ficheoutils}
	
	\subsubsection{Intervalle de Confiance (IC)}
	
	\begin{ficheoutils}{Intervalle de Confiance (IC)}
		L'\textbf{intervalle de confiance} au niveau de 95~\% pour une proportion $p$
		inconnue, calculé à partir d'un échantillon de taille $n$ et de fréquence
		observée $f$, est :
		\begin{formule}
			\[ \IC{} = \left[f - \dfrac{1}{\sqrt{n}}\ ;\ f + \dfrac{1}{\sqrt{n}}\right] \]
		\end{formule}
		\textbf{Usage :} estimer la proportion $p$ dans la population (sondages,
		enquêtes d'opinion).
		\begin{notebox}
			\textbf{Interprétation :} on dit qu'on a une confiance de 95~\% que la vraie
			proportion $p$ appartienne à cet intervalle. Plus précisément : \emph{95~\% des
				échantillons} de taille $n$ produiraient un intervalle contenant $p$.
			Condition d'application : $n \geq 30$.
		\end{notebox}
		\textbf{Marge d'erreur :} la demi-largeur de l'IC vaut environ
		$\dfrac{1}{\sqrt{n}}$. Elle diminue quand $n$ augmente.
	\end{ficheoutils}
	
	\begin{ficheoutils}{IC — Quand l'utiliser ?}
		On utilise l'\textbf{intervalle de confiance} quand la proportion $p$ de la
		population est \textbf{inconnue} et qu'on dispose d'un échantillon avec une
		fréquence observée $f$.
		
		\emph{Exemples de situations :}
		\begin{itemize}
			\item estimer la proportion de Français qui aiment le fromage (sondage sur
			1\,000 personnes) ;
			\item estimer l'intention de vote pour un candidat ;
			\item estimer la proportion de pièces défectueuses dans un lot.
		\end{itemize}
		On part de l'\textbf{observation} (échantillon) et on remonte vers la
		\textbf{population}.
	\end{ficheoutils}
	
	\begin{ficheoutils}{IC — Pourquoi l'utiliser ?}
		L'IC sert à \textbf{estimer} une proportion inconnue $p$ dans la population.
		
		Au lieu de donner une valeur unique (par exemple $f = 0{,}32$), on donne une
		\textbf{fourchette}
		$\left[f - \dfrac{1}{\sqrt{n}}\ ;\ f + \dfrac{1}{\sqrt{n}}\right]$ qui contient
		$p$ avec une confiance de 95~\%.
		
		C'est la base des sondages : la fameuse « marge d'erreur » annoncée (par exemple
		« $\pm$ 2,5 points ») correspond à $\dfrac{1}{\sqrt{n}}$ pour
		$n \approx 1\,600$.
	\end{ficheoutils}
	
	\begin{ficheoutils}{IC — Interprétation}
		Soit $\IC{} = \left[f - \dfrac{1}{\sqrt{n}}\ ;\ f + \dfrac{1}{\sqrt{n}}\right]$
		calculé sur un échantillon de taille $n$. On dit :
		\begin{itemize}
			\item « On est confiant à 95~\% que la vraie proportion $p$ de la population
			se trouve dans $\IC{}$. »
			\item ou : « Si l'on répétait un grand nombre de fois l'échantillonnage,
			environ 95~\% des intervalles calculés contiendraient $p$. »
		\end{itemize}
		\begin{notebox}
			\textbf{Piège classique :} on ne dit pas « $p$ a 95~\% de chance d'être dans
			$\IC{}$ ». La proportion $p$ est une valeur \emph{fixe} (non aléatoire) : c'est
			l'intervalle, calculé à partir d'un échantillon aléatoire, qui varie. Une fois
			calculé, soit $p$ y est, soit il n'y est pas — mais on ne peut pas le savoir.
		\end{notebox}
	\end{ficheoutils}
	
	\begin{ficheoutils}{Marge d'erreur}
		La \textbf{marge d'erreur} est la demi-largeur de l'intervalle de fluctuation ou
		de confiance. Elle vaut approximativement :
		\begin{formule}
			\[ \text{marge} \approx \dfrac{1}{\sqrt{n}} \]
		\end{formule}
		Elle \textbf{diminue} quand la taille $n$ de l'échantillon \textbf{augmente} :
		plus l'échantillon est grand, plus l'estimation est précise.
		
		\emph{Exemples :}
		\begin{itemize}
			\item $n = 100$ : marge $\approx 0{,}10$ (10 points) ;
			\item $n = 1\,000$ : marge $\approx 0{,}032$ (environ 3 points) ;
			\item $n = 2\,500$ : marge $\approx 0{,}020$ (2 points).
		\end{itemize}
		C'est pour cela que les instituts de sondage interrogent environ 1\,000 à
		2\,500 personnes : au-delà, le gain de précision devient marginal.
	\end{ficheoutils}
	
	\subsubsection{Les 4 pièges à éviter (pour ne plus jamais les confondre)}
	
	\begin{enumerate}[leftmargin=*]
		\item \textbf{Sens de lecture :} l'IF va de la \emph{Population} vers l'
		\emph{Échantillon}. L'IC va de l'\emph{Échantillon} vers la
		\emph{Population}.
		\item \textbf{Ce que je connais au départ :} pour l'IF, je connais $p$ ;
		pour l'IC, je connais $f$.
		\item \textbf{Ce que je cherche :} l'IF sert à \textbf{valider / rejeter} ;
		l'IC sert à \textbf{estimer}.
		\item \textbf{Formules (approximatives) :} IF centrée sur $p$ :
		$\left[p - \dfrac{1}{\sqrt{n}}\ ;\ p + \dfrac{1}{\sqrt{n}}\right]$ ;
		IC centrée sur $f$ :
		$\left[f - \dfrac{1}{\sqrt{n}}\ ;\ f + \dfrac{1}{\sqrt{n}}\right]$.
		Ici $n$ désigne toujours la \textbf{taille de l'échantillon}, jamais un
		effectif de modalité.
	\end{enumerate}
	
	\subsection{Vocabulaire des probabilités — ce que je modélise}
	
	Les probabilités modélisent l'incertitude. L'exemple filé : l'ouverture d'une
	\emph{loot box}.
	
	\begin{ficheprobas}{Branche probabilités (vue d'ensemble)}
		Les probabilités modélisent le hasard. Elles reposent sur trois objets
		fondamentaux :
		\begin{itemize}
			\item l'\textbf{expérience aléatoire} ;
			\item l'\textbf{univers} $\Omega$ (toutes les issues possibles) ;
			\item une \textbf{probabilité} $P$ sur $\Omega$, qui associe à chaque
			événement $A$ un nombre entre 0 et 1.
		\end{itemize}
		On y ajoute les opérations sur les événements (réunion, intersection, contraire)
		et la probabilité conditionnelle $P_B\left(A\right)$.
	\end{ficheprobas}
	
	\begin{ficheprobas}{Loi de probabilité}
		La \textbf{loi de probabilité} est la répartition des probabilités sur
		$\Omega$ : la « recette » du hasard. La somme des probabilités des issues vaut
		1.
	\end{ficheprobas}
	
	\begin{ficheprobas}{Expérience aléatoire}
		Une \textbf{expérience aléatoire} est une expérience dont on connaît toutes les
		issues possibles, mais dont on ne peut pas prévoir le résultat avant de la
		réaliser.
		
		\emph{Exemples :}
		\begin{itemize}
			\item lancer un dé à six faces ;
			\item ouvrir une loot box ;
			\item tirer un individu au sort dans une population.
		\end{itemize}
		L'ensemble des issues possibles forme l'\textbf{univers} $\Omega$.
	\end{ficheprobas}
	
	\begin{ficheprobas}{Univers $\left(\Omega\right)$}
		L'\textbf{univers}, noté $\Omega$ (oméga), est l'ensemble de toutes les issues
		possibles d'une expérience aléatoire.
		
		\emph{Exemple :} pour un dé à six faces,
		$\Omega = \left\{1\ ;\ 2\ ;\ 3\ ;\ 4\ ;\ 5\ ;\ 6\right\}$.
		
		Une \textbf{issue} est un élément de $\Omega$ (ex. : « obtenir 3 »). Un
		\textbf{événement} est un sous-ensemble de $\Omega$.
	\end{ficheprobas}
	
	\begin{ficheprobas}{Probabilité $P\left(A\right)$}
		Une \textbf{probabilité} est une fonction $P$ qui associe à chaque événement
		$A \subset \Omega$ un nombre réel dans $\left[0\ ;\ 1\right]$ mesurant la chance
		que $A$ se réalise.
		
		Propriétés fondamentales :
		\begin{itemize}
			\item $P\left(\Omega\right) = 1$ (événement certain) ;
			\item $P\left(\varnothing\right) = 0$ (événement impossible) ;
			\item la somme des probabilités des issues de $\Omega$ vaut 1.
		\end{itemize}
		Pour deux événements $A$ et $B$ :
		\begin{formule}
			\[ P\left(A \cup B\right) = P\left(A\right) + P\left(B\right) - P\left(A \cap B\right) \]
		\end{formule}
	\end{ficheprobas}
	
	\begin{ficheprobas}{Événement / Événement élémentaire}
		Un \textbf{événement} est un sous-ensemble de l'univers $\Omega$. Il est
		réalisé si l'issue obtenue appartient à cet ensemble.
		
		Un \textbf{événement élémentaire} est un événement réduit à une seule issue
		(un singleton).
		
		\emph{Exemples} avec un dé :
		\begin{itemize}
			\item « obtenir un nombre pair » $= \left\{2\ ;\ 4\ ;\ 6\right\}$ est un
			événement ;
			\item « obtenir 3 » $= \left\{3\right\}$ est un événement élémentaire.
		\end{itemize}
	\end{ficheprobas}
	
	\begin{ficheprobas}{Réunion $A \cup B$ / Intersection $A \cap B$}
		Pour deux événements $A$ et $B$ de $\Omega$ :
		\begin{itemize}
			\item la \textbf{réunion} $A \cup B$ (« $A$ ou $B$ ») est l'ensemble des
			issues appartenant à $A$, à $B$ ou aux deux ;
			\item l'\textbf{intersection} $A \cap B$ (« $A$ et $B$ ») est l'ensemble des
			issues appartenant à la fois à $A$ et à $B$.
		\end{itemize}
		Probabilités :
		\begin{formule}
			\[ P\left(A \cup B\right) = P\left(A\right) + P\left(B\right) - P\left(A \cap B\right) \]
		\end{formule}
		Si $A$ et $B$ sont \textbf{incompatibles}
		($A \cap B = \varnothing$), alors
		$P\left(A \cup B\right) = P\left(A\right) + P\left(B\right)$.
	\end{ficheprobas}
	
	\begin{ficheprobas}{Événement contraire $\bar{A}$}
		L'\textbf{événement contraire} de $A$, noté $\bar{A}$, est l'ensemble des
		issues de $\Omega$ qui \emph{ne} sont \emph{pas} dans $A$. Il se réalise
		exactement quand $A$ échoue.
		\begin{formule}
			\[ P\left(\bar{A}\right) = 1 - P\left(A\right) \]
		\end{formule}
		\emph{Exemple :} si $A$ = « obtenir un nombre pair » avec un dé
		($P\left(A\right) = 0{,}5$), alors $\bar{A}$ = « obtenir un nombre impair » et
		$P\left(\bar{A}\right) = 1 - 0{,}5 = 0{,}5$.
		
		Astuce : quand un événement est compliqué à compter directement, on calcule
		souvent la probabilité de son contraire, puis on soustrait.
	\end{ficheprobas}
	
	\begin{ficheprobas}{Probabilité conditionnelle $P_B\left(A\right)$}
		La \textbf{probabilité conditionnelle} de $A$ sachant $B$, notée
		$P_B\left(A\right)$ ou $P\left(A \mid B\right)$, mesure la probabilité que $A$
		se réalise \emph{sachant que} $B$ est déjà réalisé.
		\begin{formule}
			\[ P_B\left(A\right) = \dfrac{P\left(A \cap B\right)}{P\left(B\right)}
			\qquad \text{(si } P\left(B\right) > 0\text{)} \]
		\end{formule}
		Autrement dit, on restreint l'univers à $B$ : on ne considère que les issues où
		$B$ est réalisé, et on regarde la part de celles où $A$ se réalise aussi.
		
		\emph{Exemple :} on sait que le joueur a choisi l'Alliance (événement $B$).
		Quelle est la probabilité qu'il soit niveau maximum (événement $A$) ?
		C'est $P_B\left(A\right)$.
		
		\textbf{Formule des probabilités composées :}
		$P\left(A \cap B\right) = P_B\left(A\right) \cdot P\left(B\right)$, souvent
		utilisée pour calculer la probabilité d'une intersection.
	\end{ficheprobas}
	
	\subsubsection{Illustrations : opérations sur les événements}
	
	\begin{figure}[htbp]
		\centering
		\begin{tikzpicture}[scale=0.9, line width=0.9pt]
			% ---- Réunion ----
			\begin{scope}
				\fill[spBlue!35] (0,0) circle (1.15);
				\fill[spBlue!35] (1.7,0) circle (1.15);
				\draw (0,0) circle (1.15) node[below=1.2cm] {$A$};
				\draw (1.7,0) circle (1.15) node[below=1.2cm] {$B$};
			\end{scope}
			\node[align=center, font=\small\sffamily] at (0.85,-2.1)
			{Réunion $\left(A \cup B\right)$\\ « $A$ ou $B$ »};
		\end{tikzpicture}
		\hspace{1.2cm}
		\begin{tikzpicture}[scale=0.9, line width=0.9pt]
			% ---- Intersection ----
			\begin{scope}
				\clip (1.7,0) circle (1.15);
				\fill[spBlue!35] (0,0) circle (1.15);
			\end{scope}
			\draw (0,0) circle (1.15) node[below=1.2cm] {$A$};
			\draw (1.7,0) circle (1.15) node[below=1.2cm] {$B$};
			\node[align=center, font=\small\sffamily] at (0.85,-2.1)
			{Intersection $\left(A \cap B\right)$\\ « $A$ et $B$ »};
		\end{tikzpicture}
		\hspace{1.2cm}
		\begin{tikzpicture}[scale=0.9, line width=0.9pt]
			% ---- Contraire ----
			\fill[spBlue!35, even odd rule]
			(-1.55,-1.3) rectangle (3.25,1.3)
			(0,0) circle (1.15);
			\draw (-1.55,-1.3) rectangle (3.25,1.3);
			\draw (0,0) circle (1.15) node {$A$};
			\draw[gray, dashed] (1.7,0) circle (1.15) node[gray] {$B$};
			\node[align=center, font=\small\sffamily] at (0.85,-2.1)
			{Contraire $\left(\bar{A}\right)$\\ « non $A$ »};
		\end{tikzpicture}
		\caption{Les trois opérations sur les événements : réunion (zone colorée =
			au moins l'un des deux), intersection (seule la zone commune est colorée),
			contraire (tout sauf $A$ est coloré).}
	\end{figure}
	
	% =====================================================================
	\section{Tableaux de synthèse}
	% =====================================================================
	
	Tout le cours en un coup d'œil : nomenclature, duel IF/IC, statistiques
	descriptives, probabilités et marge d'erreur.
	
	\subsection{Nomenclature \& notations}
	
	\begin{center}
		\sffamily\bfseries Boîte à notations — qui est qui ?\\[6pt]
		\begin{tabularx}{\textwidth}{>{\raggedright\arraybackslash}p{4.6cm} X >{\raggedright\arraybackslash}p{3.2cm}}
			\toprule
			\textbf{Notion} & \textbf{Définition} & \textbf{Notation} \\
			\midrule
			Population & Ensemble de tous les individus étudiés & $\mathcal{P}$ \\
			Individu & Un élément de la population & — \\
			Échantillon & Sous-ensemble de la population & taille $n$ \\
			Taille de la population & Nombre total d'individus dans la population & $N$ \\
			Taille de l'échantillon & Nombre total d'individus prélevés & $n$ \\
			Effectif de la modalité $A$ (population) &
			Nombre d'individus de la population portant $A$ & $N_A$ \\
			Effectif de la modalité $A$ (échantillon) &
			Nombre d'individus de l'échantillon portant $A$ & $n_A$ \\
			Proportion (population) &
			Effectif de la modalité $\div$ taille de la population &
			$p = \dfrac{N_A}{N}$ \\
			Fréquence (échantillon) &
			Effectif de la modalité $\div$ taille de l'échantillon &
			$f = \dfrac{n_A}{n}$ \\
			Marge d'erreur &
			Demi-largeur de l'intervalle ; diminue quand $n$ augmente &
			$\approx \dfrac{1}{\sqrt{n}}$ \\
			\bottomrule
		\end{tabularx}
	\end{center}
	
	\subsection{Le duel : Intervalle de Fluctuation (IF) vs Intervalle de Confiance (IC)}
	
	\begin{center}
		\small
		\begin{tabularx}{\textwidth}{>{\raggedright\arraybackslash}p{3.1cm} >{\raggedright\arraybackslash}X >{\raggedright\arraybackslash}X}
			\toprule
			\textbf{Critère} & \textbf{Intervalle de Fluctuation (IF)} &
			\textbf{Intervalle de Confiance (IC)} \\
			\midrule
			Quand l'utiliser ? &
			On \textbf{connaît} la proportion théorique $p$ de la population. &
			On \textbf{ne connaît pas} $p$. On a seulement un échantillon. \\
			Ce que l'on connaît &
			La population (ex. : taux de pile d'une pièce truquée $p = 0{,}3$). &
			L'échantillon (ex. : sur 100 lancers, 40 faces, $f = 0{,}4$). \\
			Ce que l'on cherche &
			À vérifier si l'échantillon est normal… ou suspect. &
			À estimer la vraie proportion $p$ de la population. \\
			Pourquoi l'utiliser ? &
			\textbf{Prendre une décision} (valider / rejeter une hypothèse). &
			\textbf{Estimer} une réalité inconnue (sondage). \\
			Sens de lecture &
			Population $\rightarrow$ Échantillon (de la théorie vers l'observation). &
			Échantillon $\rightarrow$ Population (de l'observation vers la théorie). \\
			Centre de l'intervalle &
			Centré sur $p$ (la théorie). &
			Centré sur $f$ (l'observation). \\
			Interprétation &
			Si $f \in \IF{}$ : échantillon normal (écart dû au hasard).
			Si $f \notin \IF{}$ : on rejette l'hypothèse. &
			On est confiant à 95~\% que $p$ se trouve dans l'intervalle. \\
			Formule (approx.) &
			$\left[p - \dfrac{1}{\sqrt{n}}\ ;\ p + \dfrac{1}{\sqrt{n}}\right]$ &
			$\left[f - \dfrac{1}{\sqrt{n}}\ ;\ f + \dfrac{1}{\sqrt{n}}\right]$ \\
			Conditions d'application &
			$n \geq 25$ et $0{,}2 \leq p \leq 0{,}8$ ; seuil 95~\% (risque 5~\%). &
			$n \geq 30$ et $0{,}2 \leq f \leq 0{,}8$ ; confiance 95~\%. \\
			Exemple mental &
			« Ce dé est-il truqué ? » &
			« Combien de Français aiment le fromage ? » \\
			\bottomrule
		\end{tabularx}
	\end{center}
	
	\subsection{Statistiques descriptives — récapitulatif}
	
	\begin{center}
		\sffamily\bfseries Population contre échantillon : qui observe quoi ?\\[6pt]
		\begin{tabularx}{0.9\textwidth}{>{\raggedright\arraybackslash}p{4.6cm} >{\raggedright\arraybackslash}X >{\raggedright\arraybackslash}X}
			\toprule
			\textbf{Objet} & \textbf{Dans la population} & \textbf{Dans l'échantillon} \\
			\midrule
			Taille & $N$ & $n$ \\
			Effectif de la modalité $A$ & $N_A$ & $n_A$ \\
			Part de la modalité $A$ &
			$p = \dfrac{N_A}{N}$ (proportion théorique) &
			$f = \dfrac{n_A}{n}$ (fréquence observée) \\
			Statut & Modèle (probabilités) & Observation (statistiques) \\
			Exemple (forum) &
			$N = 10\,000$, $N_A = 3\,000$, $p = 0{,}3$ &
			$n = 500$, $n_A = 160$, $f = 0{,}32$ \\
			\bottomrule
		\end{tabularx}
	\end{center}
	
	\subsection{Probabilités — récapitulatif}
	
	\begin{center}
		\small
		\begin{tabularx}{\textwidth}{>{\raggedright\arraybackslash}p{4.2cm} X >{\raggedright\arraybackslash}p{4.6cm}}
			\toprule
			\textbf{Notion} & \textbf{Définition} & \textbf{Notation} \\
			\midrule
			Expérience aléatoire & Expérience aux résultats imprévisibles, issues connues & — \\
			Univers & Ensemble des issues possibles & $\Omega$ \\
			Probabilité & Fonction à valeurs dans $\left[0\ ;\ 1\right]$ & $P\left(A\right)$ \\
			Loi de probabilité & Répartition des probabilités sur $\Omega$ & — \\
			Événement & Sous-ensemble de $\Omega$ & $A, B$ \\
			Événement contraire & « Non $A$ », réalisé quand $A$ échoue &
			$\bar{A}$, \newline $P\left(\bar{A}\right) = 1 - P\left(A\right)$ \\
			Réunion / Intersection & $A$ ou $B$ / $A$ et $B$ & $A \cup B$ / $A \cap B$ \\
			Probabilité conditionnelle & Probabilité de $A$ sachant $B$ &
			$P_B\left(A\right)$ ou $P\left(A \mid B\right)$ \\
			\bottomrule
		\end{tabularx}
	\end{center}
	
	\subsection{Marge d'erreur : la précision coûte des individus}
	
	\begin{center}
		\sffamily\bfseries Plus l'échantillon est grand, plus la fourchette est
		resserrée — mais le gain s'essouffle.\\[6pt]
		\begin{tabular}{llll}
			\toprule
			\textbf{Taille $n$} & \textbf{Marge $\approx \dfrac{1}{\sqrt{n}}$} &
			\textbf{Lecture} & \textbf{Sondage équivalent} \\
			\midrule
			$100$ & $0{,}10$ & $\pm$ 10 points & micro-trottoir \\
			$1\,000$ & $\approx 0{,}032$ & $\pm$ 3 points & sondage standard \\
			$2\,500$ & $0{,}020$ & $\pm$ 2 points & sondage premium \\
			$10\,000$ & $0{,}010$ & $\pm$ 1 point & précision rarement rentable \\
			\bottomrule
		\end{tabular}
	\end{center}
	
	C'est pour cela que les instituts interrogent environ 1\,000 à 2\,000
	personnes : au-delà, multiplier $n$ par 4 ne divise la marge que par 2.
	
	% =====================================================================
	\section{Carte mentale}
	% =====================================================================
	
	Carte radiale : statistiques à gauche, probabilités à droite,
	échantillonnage au centre du bas. Chaque étiquette renvoie à la fiche
	détaillée correspondante du volet 1 ; les arêtes relient chaque notion à sa
	branche d'origine.
	
	\begin{figure}[htbp]
		\centering
		\resizebox{\textwidth}{!}{%
			\begin{tikzpicture}[line width=0.7pt]
				
				% ----- arêtes (dessinées d'abord) -----
				% Tronc : racine -> hubs
				\draw[spBlue!65, line width=1.1pt] (8,6.9) -- (3.7,11.6);
				\draw[spPurple!70, line width=1.1pt] (8,6.9) -- (12.3,11.6);
				\draw[spGreen!75, line width=1.1pt] (8,6.9) -- (8,4.6);
				
				% Branche statistiques : hub -> feuilles (colonne à gauche, x = 1.4)
				\draw[spBlue!55] (3.7,11.6) -- (1.4,13.3);
				\draw[spBlue!55] (3.7,11.6) -- (1.4,12.2);
				\draw[spBlue!55] (3.7,11.6) -- (1.4,11.1);
				\draw[spBlue!55] (3.7,11.6) -- (1.4,10.0);
				\draw[spBlue!55] (3.7,11.6) -- (1.4,8.9);
				% Branche probabilités : hub -> feuilles (colonne à droite, x = 14.6)
				\draw[spPurple!55] (12.3,11.6) -- (14.6,13.3);
				\draw[spPurple!55] (12.3,11.6) -- (14.6,12.2);
				\draw[spPurple!55] (12.3,11.6) -- (14.6,11.1);
				\draw[spPurple!55] (12.3,11.6) -- (14.6,10.0);
				\draw[spPurple!55] (12.3,11.6) -- (14.6,8.9);
				% Branche échantillonnage : hub -> IF / IC -> sous-nœuds
				\draw[spOrange!70, line width=1pt] (8,4.6) -- (3.6,2.6);
				\draw[spGreen!80, line width=1pt] (8,4.6) -- (12.4,2.6);
				\draw[spOrange!55] (3.6,2.6) -- (1.2,3.8);
				\draw[spOrange!55] (3.6,2.6) -- (1.2,2.5);
				\draw[spOrange!55] (3.6,2.6) -- (1.2,1.2);
				\draw[spGreen!55] (12.4,2.6) -- (14.8,3.8);
				\draw[spGreen!55] (12.4,2.6) -- (14.8,2.5);
				\draw[spGreen!55] (12.4,2.6) -- (14.8,1.2);
				
				% ----- nœud central -----
				\node[ndRoot] at (8,6.9) {Synthèse\\ Statistiques \&\\ Probabilités};
				
				% ----- branche statistiques -----
				
				\node[ndStatH] at (3.7,11.6) {Statistiques};
				\node[ndStat] at (1.4,13.3) {Population\\ $\left(N,\ p\right)$};
				\node[ndStat] at (1.4,12.2) {Individu};
				\node[ndStat] at (1.4,11.1) {Échantillon\\ $\left(n\right)$};
				\node[ndStat] at (1.4,10.0) {Caractère\\ / Modalité};
				\node[ndStat] at (1.4,8.9) {Fréquence\\ $f = \dfrac{n_A}{n}$};
				
				% ----- branche probabilités -----
				
				\node[ndProbH] at (12.3,11.6) {Probabilités};
				\node[ndProb] at (14.6,13.3) {Expérience\\ aléatoire};
				\node[ndProb] at (14.6,12.2) {Univers\\ $\left(\Omega\right)$};
				\node[ndProb] at (14.6,11.1) {Probabilité\\ $P\left(A\right)$};
				\node[ndProb] at (14.6,10.0) {Événement /\\ Élémentaire};
				\node[ndProb] at (14.6,8.9) {Réunion /\\ Intersection};
				
				% ----- branche échantillonnage -----
				\node[ndSampH] at (8,4.6) {Échantillonnage};
				\node[ndIfH] at (3.6,2.6) {Intervalle de\\ Fluctuation (IF)};
				\node[ndIf] at (1.2,3.8) {Quand ?\\ On connaît $p$.};
				\node[ndIf] at (1.2,2.5) {But ?\\ Vérifier / Décider.};
				\node[ndIf] at (1.2,1.2) {Interprétation :\\ $f \in \IF{}$ ? Normal.};
				\node[ndIcH] at (12.4,2.6) {Intervalle de\\ Confiance (IC)};
				\node[ndIc] at (14.8,3.8) {Quand ?\\ On connaît $f$.};
				\node[ndIc] at (14.8,2.5) {But ?\\ Estimer $p$.};
				\node[ndIc] at (14.8,1.2) {Confiance 95~\%\\ que $p$ y soit.};
				
			\end{tikzpicture}%
		}% fin resizebox
		\caption{Carte mentale du chapitre.}
		\label{fig:mindmap}
	\end{figure}
	
	\noindent\textbf{Légende :}
	{\color{spBlueDark}$\blacksquare$} Statistiques descriptives \quad
	{\color{spPurple}$\blacksquare$} Probabilités \quad
	{\color{spGreenDark}$\blacksquare$} Échantillonnage (IF / IC)
	
	% =====================================================================
	\section{Graphe de connaissances}
	% =====================================================================
	
	Le « chassé-croisé » de l'échantillonnage, vu comme un réseau de dépendances :
	
	\begin{itemize}
		\item l'\textbf{Intervalle de Fluctuation (IF)} part de la \textbf{Population}
		(on connaît $p$) vers l'\textbf{Échantillon} (on observe $f$) pour
		\textbf{vérifier} une hypothèse ;
		\item l'\textbf{Intervalle de Confiance (IC)} part de l'\textbf{Échantillon}
		(on connaît $f$) vers la \textbf{Population} (on cherche $p$) pour
		\textbf{estimer} une réalité.
	\end{itemize}
	
	\begin{figure}[htbp]
		\centering
		\resizebox{\textwidth}{!}{%
			\begin{tikzpicture}[
				>={Stealth[length=2.4mm]},
				line width=0.7pt]
				
				% ----- nœuds -----
				\node[gstat, text width=2.1cm] (pop)  at (7,11.8)  {Population\\ $\left(N,\ p\right)$};
				\node[gstat] (ind)  at (1.8,11.8)  {Individu};
				\node[gstat, text width=1.7cm] (ech)  at (7,9.9)  {Échantillon\\ $\left(n\right)$};
				\node[gstat, text width=2.2cm] (freq) at (11.5,9.9) {Fréquence observée\\ $f = \dfrac{n_A}{n}$};
				\node[gstat, text width=1.8cm] (car)  at (1.8,9.9)  {Caractère\\ / Modalité};
				
				\node[gtool, text width=2.2cm] (if)  at (3.3,7.4)  {\textbf{Intervalle de Fluctuation (IF)}\\ centré sur $p$};
				\node[gtool, text width=1.7cm] (marge) at (7.5,7.4) {Marge d'erreur};
				\node[gtool, text width=2.2cm] (ic)  at (11.7,7.4) {\textbf{Intervalle de Confiance (IC)}\\ centré sur $f$};
				
				\node[gprob, text width=1.9cm] (exp)  at (7,4.9)  {Expérience aléatoire};
				\node[gprob] (uni)  at (7,3.5)  {Univers $\left(\Omega\right)$};
				\node[gprob] (evt) at (7,2.1)  {Événement};
				\node[gprob, text width=1.7cm] (probA) at (3.3,2.1) {Probabilité\\ $P\left(A\right)$};
				\node[gprob, text width=1.9cm] (cond) at (11.7,2.1) {Probabilité\\ conditionnelle};
				
				% ----- arêtes statistiques (pleines, bleues) -----
				\draw[->, spBlueDark] (ind) -- node[glabel, above] {compose} (ech);
				\draw[->, spBlueDark] (ech) -- node[glabel, above] {étudie} (car);
				\draw[->, spBlueDark] (ech) -- node[glabel, above, pos=0.55] {calcule} (freq);
				\draw[->, spBlueDark] (pop) -- node[glabel, right, pos=0.6] {extrait} (ech);
				
				% ----- arêtes échantillonnage (pointillées) -----
				\draw[->, dashed, spOrange, line width=0.9pt]
				(pop) to[bend right=15] node[glabel, pos=0.3, above] {Vérifier} (if);
				\draw[->, dashed, spOrange, line width=0.9pt]
				(if) to[bend left=15] node[glabel, pos=0.5, below] {Décider} (ech);
				\draw[->, dashed, spGreenDark, line width=0.9pt]
				(freq) to[bend left=15] node[glabel, pos=0.5, right] {Estimer} (ic);
				\draw[->, dashed, spGreenDark, line width=0.9pt]
				(ic) to[bend left=12] node[glabel, pos=0.4, left] {Se fier à 95~\%} (pop);
				\draw[->, dashed, spOrange, line width=0.9pt]
				(if) -- node[glabel, above] {définit} (marge);
				\draw[->, dashed, spGreenDark, line width=0.9pt]
				(ic) -- node[glabel, above] {définit} (marge);
				
				% ----- arêtes probabilités (pleines, violettes) -----
				\draw[->, spPurple] (exp) -- node[glabel, right]{possède} (uni);
				\draw[->, spPurple] (uni) -- node[glabel, right]{contient} (evt);
				\draw[->, spPurple] (uni) -- node[glabel, above]{muni de} (probA);
				\draw[->, spPurple] (evt) -- node[glabel, above]{sachant} (cond);
				
			\end{tikzpicture}%
		}% fin resizebox
		\caption{Graphe de connaissances : réseau de dépendances entre les notions.}
		\label{fig:graphe}
	\end{figure}
	
	\noindent\textbf{Légende :}
	{\color{spBlueDark}$\blacksquare$} Statistiques descriptives \quad
	{\color{spGreenDark}$\blacksquare$} Échantillonnage (IF / IC) \quad
	{\color{spPurple}$\blacksquare$} Probabilités
	
	% =====================================================================
	\section{Simulateur IF / IC — retranscrit en exemples commentés}
	% =====================================================================
	
	Dans la version interactive, trois curseurs ($p$, $f$ et $n$) déforment les
	intervalles en direct sur une droite graduée : le meilleur moyen de
	\emph{ressentir} la marge d'erreur $\dfrac{1}{\sqrt{n}}$. Le simulateur applique
	les formules approchées du programme :
	\[
	\IF{} = \left[p - \dfrac{1}{\sqrt{n}}\ ;\ p + \dfrac{1}{\sqrt{n}}\right]
	\qquad\qquad
	\IC{} = \left[f - \dfrac{1}{\sqrt{n}}\ ;\ f + \dfrac{1}{\sqrt{n}}\right].
	\]
	Cette section en reprend les deux modes et les valeurs par défaut, sous forme
	d'exemples entièrement calculés.
	
	\subsection{Mode 1 — Tester une hypothèse (Intervalle de Fluctuation)}
	
	\textbf{Données du test :} proportion théorique $p$, taille de l'échantillon
	$n$, fréquence observée $f$.
	
	\begin{notebox}
		\textbf{Conditions d'application affichées par le simulateur :}
		$n \geq 25$ \quad et \quad $0{,}2 \leq p \leq 0{,}8$.
		Tant que ces conditions sont remplies, l'encadrement asymptotique est valable
		au seuil de 95~\% (risque d'erreur de 5~\%).
	\end{notebox}
	
	\begin{ficheoutils}{Exemple commenté : la pièce est-elle truquée ?}
		Une pièce \textbf{supposée équilibrée} ($p = 0{,}50$) est lancée
		$n = 100$ fois : on observe $f = 0{,}62$ de piles. La pièce est-elle truquée ?
		
		\medskip
		\textbf{Étape 1 — Conditions.} $n = 100 \geq 25$ \checkmark{} et
		$0{,}2 \leq 0{,}5 \leq 0{,}8$ \checkmark{} : on peut utiliser l'IF.
		
		\textbf{Étape 2 — Calcul.}
		\[
		\IF{} = \left[0{,}50 - \dfrac{1}{\sqrt{100}}\ ;\ 0{,}50 + \dfrac{1}{\sqrt{100}}\right]
		= \left[0{,}50 - 0{,}10\ ;\ 0{,}50 + 0{,}10\right]
		= \left[0{,}40\ ;\ 0{,}60\right].
		\]
		
		\textbf{Étape 3 — Verdict.} $f = 0{,}62 \notin \left[0{,}40\ ;\ 0{,}60\right]$ :
		l'échantillon est \textbf{anormal}. On \textbf{rejette} l'hypothèse que la pièce
		est équilibrée, avec un risque d'erreur de 5~\%. Pièce suspecte !
		
		\textbf{Ce que montrerait le simulateur :} en rapprochant $f$ de $0{,}60$ (par
		exemple $f = 0{,}58$), la marque $f$ rentre dans la bande : l'échantillon
		redevient « normal » — la frontière du rejet est exactement le bord de l'IF.
	\end{ficheoutils}
	
	\subsection{Mode 2 — Estimer une proportion (Intervalle de Confiance)}
	
	\textbf{Données de l'estimation :} fréquence observée $f$, taille de
	l'échantillon $n$.
	
	\begin{notebox}
		\textbf{Conditions d'application affichées par le simulateur :}
		$n \geq 30$ \quad et \quad $0{,}2 \leq f \leq 0{,}8$.
	\end{notebox}
	
	\begin{ficheoutils}{Exemple commenté : combien de joueurs de l'Alliance ?}
		Dans le forum, sur un échantillon de $n = 500$ membres, $n_A = 160$ jouent
		l'Alliance, donc $f = \dfrac{160}{500} = 0{,}32$. Estimer la vraie proportion
		$p$ de l'Alliance dans le forum.
		
		\medskip
		\textbf{Étape 1 — Conditions.} $n = 500 \geq 30$ \checkmark{} et
		$0{,}2 \leq 0{,}32 \leq 0{,}8$ \checkmark{} : on peut utiliser l'IC.
		
		\textbf{Étape 2 — Marge.} $\dfrac{1}{\sqrt{500}} \approx 0{,}045$.
		
		\textbf{Étape 3 — Intervalle.}
		\[
		\IC{} = \left[0{,}32 - 0{,}045\ ;\ 0{,}32 + 0{,}045\right]
		\approx \left[0{,}275\ ;\ 0{,}365\right].
		\]
		
		\textbf{Verdict.} On est \textbf{confiant à 95~\%} que la vraie proportion $p$
		de l'Alliance dans le forum se trouve entre environ $0{,}28$ et $0{,}37$.
		
		\textbf{Insight du simulateur :} multiplier $n$ par 4 ne divise la marge que
		par 2 (car $\dfrac{1}{\sqrt{4n}} = \dfrac{1}{2}\cdot\dfrac{1}{\sqrt{n}}$) :
		la précision coûte cher.
	\end{ficheoutils}
	
	\subsection{À retenir du simulateur}
	
	\begin{itemize}
		\item la \textbf{largeur} de la bande colorée ne dépend que de $n$
		(marge $\approx 1/\sqrt{n}$), jamais de $p$ ni de $f$ ;
		\item l'IF est \textbf{centré sur $p$}, l'IC est \textbf{centré sur $f$} :
		glisser le curseur $p$ (ou $f$) déplace la bande entière ;
		\item la décision IF est binaire : $f \in \IF{} \Rightarrow$ normal,
		$f \notin \IF{} \Rightarrow$ rejeter au risque de 5~\% ;
		\item pour diviser la marge par 2, il faut \textbf{quadrupler} $n$.
	\end{itemize}
	
	% =====================================================================
	\section{Quiz d'auto-évaluation (10 questions)}
	% =====================================================================
	
	Dix questions couvrant la nomenclature, le duel IF/IC et les probabilités.
	Cochez la bonne réponse, puis confrontez-la au corrigé de la
	section~\ref{sec:corrige} : chaque explication renvoie à la fiche
	correspondante du volet 1.
	
	\begin{enumerate}[leftmargin=*, label=\textbf{Q\arabic*.}]
		\item Une population compte 10\,000 membres, dont 3\,000 choisissent
		l'Alliance. Que vaut la proportion $p$ ?
		\begin{enumerate}[label=(\alph*)]
			\item $0{,}3$
			\item $3\,000$
			\item $0{,}03$
			\item $3\,000 \div 500$
		\end{enumerate}
		
		\item Dans un échantillon de 500 membres, 160 choisissent l'Alliance. Que vaut
		la fréquence $f$ ?
		\begin{enumerate}[label=(\alph*)]
			\item $160$
			\item $0{,}32$
			\item $0{,}5$
			\item $500$
		\end{enumerate}
		
		\item Que désigne la lettre $n$ (sans indice) ?
		\begin{enumerate}[label=(\alph*)]
			\item Le nombre d'individus portant la modalité $A$
			\item La taille de la population
			\item La taille de l'échantillon
			\item La fréquence observée
		\end{enumerate}
		
		\item Une pièce équilibrée ($p = 0{,}5$) est lancée 100 fois : on observe
		$f = 0{,}62$ piles. Quel outil décide si elle est truquée ?
		\begin{enumerate}[label=(\alph*)]
			\item L'intervalle de confiance
			\item L'intervalle de fluctuation
			\item La probabilité conditionnelle
			\item La loi de probabilité
		\end{enumerate}
		
		\item Un institut interroge 1\,000 personnes pour estimer une proportion
		inconnue. Quel outil utilise-t-il ?
		\begin{enumerate}[label=(\alph*)]
			\item L'intervalle de fluctuation
			\item L'intervalle de confiance
			\item L'univers $\Omega$
			\item La marge d'erreur
		\end{enumerate}
		
		\item L'intervalle de fluctuation au seuil de 95~\% est centré sur… ?
		\begin{enumerate}[label=(\alph*)]
			\item $f$
			\item $p$
			\item $n$
			\item $0{,}95$
		\end{enumerate}
		
		\item Au seuil de 95~\%, on calcule l'IF et on trouve $f \notin \IF{}$.
		Conclusion ?
		\begin{enumerate}[label=(\alph*)]
			\item L'échantillon est normal
			\item On rejette l'hypothèse, avec un risque d'erreur de 5~\%
			\item $p$ est forcément hors de l'IF
			\item L'expérience n'est pas aléatoire
		\end{enumerate}
		
		\item Quelle phrase est correcte pour un IC à 95~\% ?
		\begin{enumerate}[label=(\alph*)]
			\item « $p$ a 95~\% de chances d'être dans l'intervalle »
			\item « On est confiant à 95~\% que $p$ appartient à l'intervalle »
			\item « 95~\% de la population est dans l'intervalle »
			\item « $f$ appartient à l'intervalle dans 95~\% des cas »
		\end{enumerate}
		
		\item La marge d'erreur $1/\sqrt{n}$…
		\begin{enumerate}[label=(\alph*)]
			\item augmente avec $n$
			\item diminue quand $n$ augmente
			\item ne dépend pas de $n$
			\item vaut toujours 5~\%
		\end{enumerate}
		
		\item Que vaut $P\left(A \cup B\right)$ ?
		\begin{enumerate}[label=(\alph*)]
			\item $P\left(A\right) + P\left(B\right)$
			\item $P\left(A\right) + P\left(B\right) - P\left(A \cap B\right)$
			\item $P\left(A\right) \times P\left(B\right)$
			\item $P\left(A\right) - P\left(B\right)$
		\end{enumerate}
	\end{enumerate}
	
	\subsection{Corrigé détaillé}
	\label{sec:corrige}
	
	\begin{enumerate}[leftmargin=*, label=\textbf{Q\arabic*.}]
		\item \textbf{(a) $0{,}3$.} \ $p = N_A / N = 3\,000 / 10\,000 = 0{,}3$ :
		effectif de la modalité divisé par la taille de la population.
		\emph{Fiche : Population.}
		\item \textbf{(b) $0{,}32$.} \ $f = n_A / n = 160 / 500 = 0{,}32$ : effectif de
		la modalité dans l'échantillon, divisé par la taille de l'échantillon.
		\emph{Fiche : Fréquence.}
		\item \textbf{(c) La taille de l'échantillon.} \ $n$ est la taille de
		l'échantillon. Un effectif de modalité porte toujours un indice :
		$n_A$, $n_i$, $N_A$, \ldots \emph{Fiche : Échantillon.}
		\item \textbf{(b) L'intervalle de fluctuation.} \ $p$ est connue et on teste un
		échantillon : c'est le rôle de l'IF, centré sur $p$. Ici $f = 0{,}62$ sort
		de $\left[0{,}4\ ;\ 0{,}6\right]$ : pièce suspecte !
		\emph{Fiche : Intervalle de Fluctuation.}
		\item \textbf{(b) L'intervalle de confiance.} \ $p$ est inconnue, on remonte de
		l'observation $f$ vers la population : c'est l'IC, centré sur $f$.
		\emph{Fiche : Intervalle de Confiance.}
		\item \textbf{(b) $p$.} \ L'IF est centré sur la théorie $p$. C'est l'IC qui
		est centré sur l'observation $f$. Le 95~\% est un niveau de confiance, pas
		un centre. \emph{Fiche : Intervalle de Fluctuation.}
		\item \textbf{(b) On rejette l'hypothèse, avec un risque d'erreur de 5~\%.} \
		Un écart trop grand pour être dû au hasard : on rejette l'hypothèse — en
		acceptant un risque de 5~\% de se tromper.
		\emph{Fiche : IF — Interprétation.}
		\item \textbf{(b) « On est confiant à 95~\% que $p$ appartient à
			l'intervalle ».} \ $p$ est une valeur fixe : c'est l'intervalle qui varierait
		si l'on répétait l'échantillonnage. Environ 95~\% des intervalles calculés
		contiendraient $p$. \emph{Fiche : IC — Interprétation.}
		\item \textbf{(b) Elle diminue quand $n$ augmente.} \ La marge diminue comme
		$1/\sqrt{n}$ : diviser la marge par 2 exige de multiplier $n$ par 4. La
		précision coûte cher. \emph{Fiche : Marge d'erreur.}
		\item \textbf{(b) $P\left(A\right) + P\left(B\right) - P\left(A \cap B\right)$.} \
		On soustrait l'intersection pour ne pas la compter deux fois. Si $A$ et $B$
		sont incompatibles ($A \cap B = \varnothing$), alors
		$P\left(A \cup B\right) = P\left(A\right) + P\left(B\right)$.
		\emph{Fiche : Réunion / Intersection.}
	\end{enumerate}
	
	\begin{tcolorbox}[colback=spBlue!5, colframe=spBlueDark!75, boxrule=0.8pt,
		arc=2mm, left=3mm, right=3mm, top=2mm, bottom=2mm,
		title={\sffamily\bfseries Barème de commentaires}, coltitle=white,
		colbacktitle=spBlueDark, fonttitle=\bfseries]
		\begin{itemize}
			\item \textbf{10/10} — Sans faute : le chassé-croisé IF / IC n'a plus de
			secret pour vous !
			\item \textbf{8--9} — Très solide ! Encore un ou deux pièges de notation à
			verrouiller.
			\item \textbf{6--7} — Bonne base : revoyez les fiches des questions ratées.
			\item \textbf{$\leq$ 5} — Reprenez le cours en fiches (volet 1) : chaque
			erreur renvoie à sa fiche.
		\end{itemize}
	\end{tcolorbox}
	
	% =====================================================================
	\section*{Pour aller plus loin}
	\addcontentsline{toc}{section}{Pour aller plus loin}
	% =====================================================================
	
	Terminale générale, option Mathématiques complémentaires — synthèse complète.
	Formules approchées du programme : IF et IC au seuil de confiance 95~\%.
	Franchement testable : vérifiez que $f \in \IF{}$ avant de rejeter un
	échantillon !
	
\end{document}